\section{Discussion}
\label{sec:discussion}

\textbf{Limitations and future directions.}
First, \bugs{}'s collection pipeline is Python-centric.
The mechanisms to identify programmatic objects (e.g. functions, classes) and perform transformations rely heavily on the Python specific \texttt{ast} library.
That said, \bugs{}'s collection strategy is transferable to other languages.
Second, due to both compute/budget constraints and our work's primary stance as a dataset contribution, we only include fine-tuning as a demonstration of \bugs{}'s effectiveness.
We do not explore other training techniques such as reasoning capabilities elicited via reinforcement learning.
% To this end, we hope the community finds \bugs{} helpful to support future work such as integrating reasoning and agentic capabilities for AI software engineering.

\textbf{Conclusion.}
We introduce \bugs{}, a dataset of $50$k software engineering task instances from across $128$ real world GitHub repositories.
\bugs{} collection pipeline allows us to scale up task instances, environments, and trajectories at a fraction of prior costs without sacrificing faithfulness to open source software development practices.
Using \bugs{}, we train \texttt{SWE-agent-LM-32B}, achieving a state-of-the-art $40.2$\% on SWE-bench Verified.
Our experiments show how \bugs{} can be used to identify fundamental trends about developing SWE-agents.
We believe \bugs{} provides the foundational data and infrastructure needed to train software engineering agents in a truly scalable manner.
